% SPDX-License-Identifier: Apache-2.0 OR CC-BY-4.0
\documentclass[11pt]{article}
\usepackage{amsmath,amssymb,amsthm,hyperref}
\hypersetup{hidelinks,pdftitle={Prime-factor density reversals and the spectrum of bounded prime gaps},pdfauthor={Jonas Whidden}}
\setlength{\emergencystretch}{3em}
\newtheorem{theorem}{Theorem}[section]
\newtheorem{lemma}[theorem]{Lemma}
\newtheorem{corollary}[theorem]{Corollary}
\theoremstyle{definition}
\newtheorem{definition}[theorem]{Definition}
\newcommand{\LL}{\mathcal L}
\newcommand{\EE}{\mathcal E}
\newcommand{\N}{\mathbb N}
\title{Prime-factor density reversals and the spectrum of bounded prime gaps}
\author{Jonas Whidden}
\date{Research draft, 28 September 2026}
\begin{document}
\maketitle
\begin{abstract}
We study strict descents followed by later strict ascents in the density
of the $k$th smallest distinct prime factor. For a fixed local probability
law satisfying $\eta_p^{-1}=p/\nu+\beta+o(1)$, with $\nu>0$, every rank
has a finite, attained maximal number $N_\eta(k)$ of separated reversals.
Its double-logarithmic upper growth rate is exactly
$\sup_{H\ge2}\gamma_H/(H+\nu)$, where $\gamma_H$ measures the actual
frequency of consecutive prime gaps at most $H$. The positive rate has
a finite phase profile and applies to limiting factorization densities
of products of distinct affine forms in the sum of two primes.
For the ordinary and generic odd laws we also obtain double-exponential
bounds at every sufficiently large rank and an exact equivalence between
recurring short gaps and unbounded late-ascent supply. The proofs combine
uniform symmetric-polynomial estimates with a finite selection argument
using periodic composite blocks. The maintained results are formalized
in Lean. The transfers do not establish a new numerical prime-gap bound.
\end{abstract}

\section{Introduction and historical context}\label{sec:intro}
Let $d_k(p)$ be the natural density of positive integers whose $k$th
smallest \emph{distinct} prime factor is $p$. The question whether
$p\mapsto d_k(p)$ first rises and then falls is now called Erd\H{o}s
Problem 690. Our subject is the finer structure left after unimodality
fails: the number of separated reversals, their possible locations, and
the arithmetic information encoded by their growth.

Erd\H{o}s formulated the question in 1979~\cite[p.~75]{Erdos}.
The same paper discusses the typical scale $\log\log p_k(n)\sim k$
for prime factors of integers. This concerns where probability mass
lies, not the mode of the individual-prime density. Erd\H{o}s--Tenenbaum
\cite[Corollary~5]{ET} later corrected the earlier assertion about the
mode: a maximizing prime satisfies $\log p\sim k/\log k$, rather than
$\log p\sim k$. Neither statement counts reversals. The occurrence of
a double exponential in classical prime-factor distribution does not
by itself supply the theorem below.

Cambie~\cite[Theorem~5]{Cambie} proved unimodality for $k=1,2,3$ and
verified its failure for $4\le k\le20$. His Theorem~3 gives
superpolynomially many local maxima for the different density
$\delta_1(n,m)$ in Erd\H{o}s Problem 692: the density of integers having
exactly one divisor in $\{n+1,\ldots,m-1\}$. Its moving variable is $m$,
with $n$ fixed. It is not a reversal count for $d_k(p)$.

Wang--Crapis~\cite[Theorem~1.1]{WC} proved non-unimodality for every
$k\ge4$. Their proof compares a symmetric-polynomial ratio with prime
gaps and produces a descent followed by an ascent. We extend the
comparison to local laws and make the construction quantitative.
The question here is how the maximal number of witnesses grows with
$k$. Its exact rate also requires a converse: a sufficiently late
ascent must come from an appropriately short gap.

The contribution is the resulting frequency spectrum, its realization
in prime-pair arithmetic, and its late-ascent interpretation.
Quantitative bounded-gap results are existing analytic inputs, not
new sieve results of this paper. The
\href{https://github.com/kimihiro64/prime-factor-oscillations}{public repository}
contains maintained proofs and a provenance ledger. This is an
AI-assisted research draft; independent expert review and assessment
of novelty remain outstanding.

\section{The argument in plain English}\label{sec:plain}
Choose a position in an integer's ordered list of prime factors,
say the tenth distinct prime factor. For each prime, ask what fraction
of sampled integers put that prime in the chosen position. These
fractions form a sequence. A reversal means that the sequence drops
and later rises again. We count occurrences separated far enough
that they do not reuse the same stretch of the sequence.

Two effects compete. Moving to a larger prime lowers its own chance
of dividing the integer. At the same time, the additional smaller
primes make it easier to have the required number of preceding
factors. An exact formula measures this competition. On the relevant
scale, a short gap between consecutive primes favors a rise, whereas
a long gap favors a drop.

The useful measuring scale is the logarithm of the logarithm of the
prime. On that scale the available rise threshold is approximately
the rank divided by the scale. This explains the double-exponential
locations in the argument. Eventually even the shortest possible
prime gap is too large for a rise, so every fixed rank still has a
final decreasing tail.

Many short gaps alone are not enough to count reversals. We must
put a long gap before each selected short gap and keep the witnesses
distinct. Periodic blocks of composite integers provide barriers
that consecutive primes must jump. Grouping short gaps between the
barriers and taking alternate occupied groups produces separated
witnesses, losing only a fixed factor in their number.

Finally, different short-gap sizes can have different frequencies.
Each size contributes its frequency divided by a threshold cost.
The greatest contribution is exactly the reversal growth rate.
For products of affine prime sums, the number of distinct factors
changes that cost, while finitely many exceptional primes change
only the initial details. The last sections explain the prime-gap
information this recovers and the limits of that conclusion.

\section{Definitions and main results}\label{sec:statements}
All logarithms are natural. Index the primes by $p_0=2,p_1=3,\ldots$,
and write $\LL(x)=\log\log x$ for $x>1$ and
$\EE(t)=\exp(\exp t)$. Ranks $k$ are positive integers.
A fixed \emph{local law} consists of probabilities $0\le\eta_p\le1$.
Define the finite products
\begin{equation}\label{eq:density}
 M_r(p)=[z^r]\prod_{q<p}(1-\eta_q+\eta_qz),\qquad
 D_{\eta,k}(p)=\eta_pM_{k-1}(p),
\end{equation}
with $M_r=0$ for $r<0$. For the ordinary law $\eta_p=1/p$,
the Chinese remainder theorem identifies $D_{\eta,k}$ with $d_k$.
The second canonical law is
\begin{equation}\label{eq:odd}
 \eta_2^{\mathrm{odd}}=0,\qquad
 \eta_p^{\mathrm{odd}}=\frac{p-2}{(p-1)^2}\quad(p>2).
\end{equation}

\begin{definition}\label{def:reversal}
A reversal at rank $k$ consists of indices $d<a$ with
\[
 D_{\eta,k}(p_{d+1})<D_{\eta,k}(p_d),\qquad
 D_{\eta,k}(p_a)<D_{\eta,k}(p_{a+1}).
\]
A family of $m$ such witnesses is separated if $a_i+1<d_j$ for
$i<j$. Its maximal attainable size, when finite, is $N_\eta(k)$.
\end{definition}
This counts raw-density descent--ascent witnesses, not oscillations
of a centered error or a smoothed approximation.

For integers $H\ge0$, set
\begin{equation}\label{eq:spectrum}
\begin{split}
 G_H(X)&=\#\{i:p_i\le X,\ p_{i+1}-p_i\le H\},\\
 \gamma_H&=\limsup_{\substack{X\to\infty\\X\in\N}}
       \frac{\log\log(3+G_H(X))}{\LL(X)},\\
 \Lambda(\nu)&=\sup_{H\ge2}\frac{\gamma_H}{H+\nu}\quad(\nu>0).
\end{split}
\end{equation}
The $3$ keeps the logarithms defined when counts vanish.
We have $0\le\gamma_H\le1$ and $\gamma_0=\gamma_1=0$.

\begin{theorem}[Exact reversal spectrum]\label{thm:spectrum}
Suppose $\nu>0$, $\beta\in\mathbb R$, and, along sufficiently large primes,
\begin{equation}\label{eq:smooth}
 0<\eta_p<1,\qquad
 \eta_p^{-1}=p/\nu+\beta+\rho(p),\qquad \rho(p)\longrightarrow0.
\end{equation}
Every fixed-rank sequence $D_{\eta,k}$ is eventually nonincreasing,
its reversal maximum is finite and attained, and
\begin{equation}\label{eq:rate}
 \limsup_{k\to\infty}\frac{\log\log(3+N_\eta(k))}{k}
       =\Lambda(\nu)>0.
\end{equation}
The law, its parameters and its exceptional coordinates are fixed
before $k$ tends to infinity.
\end{theorem}
The weaker estimate $\eta_p=\nu/p+O(p^{-2})$ alone does not supply
the vanishing difference of reciprocal errors used in this theorem.

\begin{theorem}[Canonical bounds]\label{thm:canonical}
For either canonical law and every $\varepsilon>0$, for all
sufficiently large $k$,
\[
 \EE(k/602)\le N_\eta(k)\le\EE((1/3+\varepsilon)k).
\]
Each fixed-rank canonical sequence is eventually strictly decreasing.
\end{theorem}
The positive constant uses a conservative quantitative width-$600$
input; it is not a claim about the best known numerical prime-gap bound.

\begin{theorem}[Affine arithmetic realization]\label{thm:affine}
Fix $m\ge1$, positive integers $c_j$, and integers $d_j$, with
$c_id_j\ne c_jd_i$ for $i\ne j$. Put
\[
 F(p,q)=\prod_{j=1}^m(c_j(p+q)+d_j).
\]
Sample all ordered prime pairs $p,q\le X$, including the diagonal,
with denominator $\pi(X)^2$. At every fixed rank $k$ and prime
$\ell$, the probability that $\ell$ is the $k$th smallest distinct
prime factor of $|F(p,q)|$ has a limit $D_{F,k}(\ell)$; zero
outputs are excluded from the event. The maximal separated reversal
count $N_F(k)$ is finite and attained, and
\[
 \limsup_{k\to\infty}\frac{\log\log(3+N_F(k))}{k}=\Lambda(m)>0.
\]
\end{theorem}
The size limit comes first, with family, rank and prime fixed.
Outputs are weighted by their prime-pair representations.

\section{The exact comparison and uniform estimates}\label{sec:ratio}
\subsection{First differences}
For consecutive primes $p<q$ with interior probabilities,
\[
 M_r(q)=(1-\eta_p)M_r(p)+\eta_pM_{r-1}(p).
\]
Direct subtraction gives
\begin{equation}\label{eq:diff}
 D_{\eta,k}(q)-D_{\eta,k}(p)=
 \eta_p\eta_q\{M_{k-2}(p)-T_\eta(p,q)M_{k-1}(p)\},
\end{equation}
where $T_\eta(p,q)=\eta_q^{-1}-\eta_p^{-1}+1$.
When $M_{k-1}(p)>0$, an ascent is equivalent to
$M_{k-2}(p)/M_{k-1}(p)>T_\eta(p,q)$.
The ordinary threshold is $q-p+1$; for the odd law at $p>2$ it is
\[
 q-p+1-\frac{q-p}{(p-2)(q-2)}.
\]

\subsection{Finite exceptions and symmetric polynomials}
Only finitely many coordinates in \eqref{eq:smooth} can be forced
($\eta_p=1$) or excluded ($\eta_p=0$). Let $f$ count the forced
coordinates. Past them, each contributes a factor $z$ to the
generating polynomial; excluded coordinates contribute $1$.
Retain all interior coordinates, including the exceptional ones,
and set
\[
 w_p=\frac{\eta_p}{1-\eta_p},\qquad
 S(p)=\sum_{\substack{q<p\\0<\eta_q<1}}w_q.
\]
For $k>f$, put $r=k-1-f$. The common positive product of
$1-\eta_q$ cancels in the ratio, leaving $E_{r-1}(p)/E_r(p)$,
where $E_r$ is the elementary symmetric polynomial of degree $r$
in the finite list of interior odds.

\begin{lemma}\label{lem:symmetric}
For at least $r\ge1$ positive weights, let $W_{r-1}$ be the largest
sum of $r-1$ distinct weights. Then
\begin{equation}\label{eq:ratio}
 \frac rS\le\frac{E_{r-1}}{E_r}\le\frac r{S-W_{r-1}},
\end{equation}
and the right denominator is positive.
\end{lemma}
\begin{proof}
The terms of $SE_{r-1}$ which choose a new index total $rE_r$.
The remaining terms are
\[
 \sum_{|A|=r-1}\left(\sum_{j\in A}w_j\right)\prod_{j\in A}w_j,
\]
between zero and $W_{r-1}E_{r-1}$. Rearrangement proves the bounds.
At least one positive weight remains after any $r-1$ are removed.
\end{proof}
This specializes to Wang--Crapis \cite[Lemmas~3.1--3.2]{WC};
the present use also retains the exact forced-rank shift.

Expanding \eqref{eq:smooth} gives $w_p=\nu/p+O(p^{-2})$.
The reciprocal-prime estimate and the convergent quadratic tail give
\begin{equation}\label{eq:mertens}
 S(p)=\nu\LL(p-1)+O_\eta(1),\qquad
 W_{r-1}=O_\eta(1+\log(r+1)).
\end{equation}
For the second estimate, dominate all but finitely many odds by
$C/p$ and use $p_j\ge j+2$. The finitely many exceptional odds
contribute a fixed constant; their order need not be monotone.

\begin{lemma}[Uniform sign margins]\label{lem:bands}
Fix $0<u\le v$. Uniformly in the band
$uk\le\LL(p-1)\le vk$, as $k\to\infty$,
\[
 \frac{E_{r-1}(p)}{E_r(p)}
       =\frac{k}{\nu\LL(p-1)}+o(1),\qquad r=k-1-f.
\]
Gaps at most $H$ give strict ascents throughout the band if
$(H+\nu)v<1$. Gaps at least $D$ give strict descents there if
$Du>2$, for sufficiently large rank and lower prime.
\end{lemma}
\begin{proof}
The prime number theorem ensures more than $r$ interior primes
below $p$ uniformly in this band. Equations \eqref{eq:ratio} and
\eqref{eq:mertens} have numerator $k+O_\eta(1)$ and denominator
$\nu\LL(p-1)+O_\eta(\log(k+1))$, bounded below by a positive
multiple of $k$. This proves the uniform ratio statement.

Writing $g=q-p$, the threshold is
$g/\nu+1+\rho(q)-\rho(p)$. Its error is bounded by twice the
supremum of $|\rho|$ beyond $p$, even for large $g$.
For $g\le H$ the limiting lower ratio is
$1/(\nu v)>(H+\nu)/\nu$.
For $g\ge D$ the upper ratio is
$1/(\nu u)+o(1)<D/\nu$. Both margins are strict.
\end{proof}
Without an upper band endpoint, the same estimates give a
one-sided bound: on $\LL(p-1)\ge bk$ the ratio is at most
$1/(\nu b)+o(1)$, uniformly for fixed $b>0$.

\subsection{Every fixed rank}
For fixed $r\ge1$, the ratio tends to zero by
\eqref{eq:ratio}--\eqref{eq:mertens}. Consecutive sufficiently large
primes have gap at least two, so their threshold is uniformly
positive. Equation \eqref{eq:diff} gives strict descent for
supported degrees. If $k\le f$, the tail is zero. If $k=f+1$,
the remaining degree is zero, its predecessor coefficient vanishes,
and the same identity again gives descent.
Thus every fixed-rank tail is nonincreasing.

Once all ascent indices lie below an integer $I$, a separated
family has size at most $I$. Attainable sizes form a nonempty
bounded subset of $\N$ (zero is attainable), so their maximum
is attained. The canonical laws have no forced coordinates
and have positive supported tails, giving strict eventual decrease.

\section{From short gaps to separated reversals}\label{sec:selection}
The following finite argument supplies the counting step and its loss.

\begin{lemma}[Periodic composite barriers]\label{lem:cells}
Fix integers $D\ge H+1$ and $M=D!$. Let $\mathcal S$ be a
nonempty finite set of integer starts such that
\[
 X\le n,\quad n+H\le Y,\quad
 \#([n,n+H]\cap\{\text{primes}\})\ge2,\quad X\ge M+2.
\]
There are $m$ separated pairs of consecutive-prime gaps, a gap at
least $D$ followed by one at most $H$, all endpoints in $[X,Y]$,
with
\begin{equation}\label{eq:cellloss}
 |\mathcal S|\le2(M+H)(m+1).
\end{equation}
\end{lemma}
\begin{proof}
Choose consecutive primes $p<q$ in each window, and assign them
to $c=\lfloor(p-2)/M\rfloor$. Since $cM+j$ is composite for
$2\le j\le D$, and $c\ge1$,
\[
 cM+D+1\le p<q\le(c+1)M+1.
\]
Indeed a prime $q$ beyond the last endpoint but within $H<D$
of $p$ would lie in the next composite block. A short pair cannot
cross that barrier.

Starts assigned to cell $c$ lie between $cM+D+1-H$ and
$(c+1)M+1$, so there are at most $M+H$ of them.
Order the occupied cells $c_0<\cdots<c_{t-1}$ and choose one
pair in each, with lower-prime index $u_j$. Between chosen
primes in two successive occupied cells, a consecutive-prime gap
crosses the destination cell's whole composite block. Its length
is at least $D$, and it precedes the chosen short gap there.

Take destination cells $c_2,c_4,\ldots,c_{2m}$, where
$m=\lfloor(t-1)/2\rfloor$. Choose each crossing after the lower
prime selected in the preceding occupied cell. The unused occupied
cell between successive destinations ensures $a_i+1<d_{i+1}$.
Every crossing lies between selected endpoints, so remains in
$[X,Y]$. Finally $t\le2(m+1)$ and
$|\mathcal S|\le(M+H)t$, proving \eqref{eq:cellloss}.
\end{proof}
No uniform distribution among cells is needed. Empty stretches
are harmless, since a crossing still exists between selected
occupied cells.

Fix $u,v>0$ with $(H+\nu)v<1$ and choose an integer
$D\ge H+1$ with $D>2/u$. If all primes from $X$ through $Y+H$
lie in the sign band, the two lemmas imply, at sufficiently large $k$,
\begin{equation}\label{eq:bandcount}
 G_H(Y)\le X+B(N_\eta(k)+1),\qquad B=2(D!+H).
\end{equation}
At most $X$ lower primes are deleted below $X$; each retained
short gap supplies a window starting at its lower prime.
The upper endpoint is $Y+H$, since a counted gap may end above $Y$.
All constants are fixed before $k$ grows.

\begin{corollary}[Positive-power transfer]\label{cor:power}
Suppose that for fixed $H$ and $\delta>0$, for every sufficiently
large $X$ at least $X^\delta$ integer starts in $[X,2X]$ have
two primes in their width-$H$ windows. For every
$0<a<1/(H+\nu)$,
\[
 N_\eta(k)\ge\EE(ak)\quad\hbox{for all sufficiently large }k.
\]
\end{corollary}
\begin{proof}
Choose $a<b<1/(H+\nu)$ and $X=\lceil\EE(bk)\rceil$.
Every endpoint in $[X,2X+H]$ has $\LL(p-1)=bk+o(1)$.
Fixed $u<b<v<1/(H+\nu)$ therefore provide the sign band.
The cell lemma gives $N_\eta(k)+1\ge X^\delta/B$.
Since $\log(X^\delta/B)=\delta\exp(bk)+O(1)$ and $a<b$,
the right side eventually exceeds $1+\EE(ak)$.
\end{proof}

We use the existing quantitative input
\begin{equation}\label{eq:clusterinput}
 \#\{n\in[X,2X]\cap\N:
       \#([n,n+600]\cap\{\text{primes}\})\ge2\}\ge X^{1/3}
\end{equation}
for all sufficiently large $X$. Its maintained declaration is
\path{PrimeGaps.eventually_many_prime_windows_600}.
Section~\ref{sec:formal} records its provenance. Quantitative
prime-cluster theory, including Maynard's work~\cite{Maynard},
is the analytic background. We use this count, not merely
infinitude of bounded gaps. With $\nu=1$ the corollary applies
at $a=1/602<1/601$.

\section{Proof of the exact spectrum}\label{sec:sharp}
\subsection{Lower bound: arbitrary frequency subsequences}
Fix $H$ and $0<a<\gamma_H/(H+\nu)$. Choose
\[
 0<g<\gamma_H,\qquad
 0<v_0<v<\frac1{H+\nu},\qquad a<gv_0,
\]
then $0<u<a$ and $a<b<gv_0$.
There are arbitrarily large integers $Y$ with
$3+G_H(Y)>\EE(g\LL(Y))$. Assign
$k=\lceil\LL(Y)/v_0\rceil$. The bounds
$v_0k-v_0<\LL(Y)\le v_0k$ imply
$3+G_H(Y)>\EE(bk)$ at arbitrarily large assigned ranks.

Take $X=\lceil\EE(uk)\rceil+1$. The primes between $X$ and
$Y+H$ lie in a band with positive lower coefficient and upper
coefficient $v$. Equation \eqref{eq:bandcount} applies.
If $N_\eta(k)<\EE(ak)$, then
\[
 3+G_H(Y)\le3+X+B(1+\EE(ak))<\EE(bk)
\]
eventually, a contradiction. Therefore $N_\eta(k)\ge\EE(ak)$
for arbitrarily large $k$. Supremizing over $H$ proves the lower
bound $\Lambda(\nu)$ in \eqref{eq:rate}.
The argument keeps the actual frequency subsequence; it does not
replace it by an all-large-$X$ hypothesis.

\subsection{Upper bound: each gap retains its own cutoff}
For fixed integer $g\ge2$ and $\varepsilon>0$, the one-sided
ratio bound implies that an ascent across a gap of size $g$
must satisfy, apart from fixed initial exceptions,
\begin{equation}\label{eq:cutoff}
 \LL(p-1)<\left(\frac1{g+\nu}+\varepsilon\right)k.
\end{equation}
The same estimate excludes all gaps greater than $J$ above
$\LL(p-1)=uk$ if $1/(J+1+\nu)<u$.
In both statements, the reciprocal error is uniformly small
beyond the lower prime.

Consequently
\begin{equation}\label{eq:envelope}
 N_\eta(k)\le T_k+\sum_{2\le g\le J}G_g(U_g(k)),
\end{equation}
where we may take
\[
 T_k=\lceil\EE(uk)\rceil+O(1),\qquad
 U_g(k)=\left\lceil\EE\left(
       \left(\frac1{g+\nu}+\varepsilon\right)k\right)\right\rceil+1.
\]
Each witness has a distinct ascent index. There are at most $T_k$
initial ones; all others have an actual gap $g\le J$ and are
counted in that gap's term. Overcounting by the cumulative
$G_g$ only strengthens the upper bound.

Let $A>\Lambda(\nu)$. Choose $0<u<A$ and then $J$ as above.
For each fixed $g$ and $\delta>0$, eventually
\[
 \log\log(3+G_g(U))\le(\gamma_g+\delta)\LL(U).
\]
Since $\gamma_g/(g+\nu)\le\Lambda(\nu)$, choose
$\varepsilon,\delta>0$ small enough that
\[
 (\gamma_g+\delta)
       \left(\frac1{g+\nu}+\varepsilon\right)<A
       \quad(2\le g\le J).
\]
There are only finitely many such comparisons. Rounding changes
their double-logarithmic cutoffs by $o(1)$. Each term of
\eqref{eq:envelope}, including $T_k$, has rate strictly below $A$.
A fixed finite sum has rate at most the greatest term rate.
Letting $A$ decrease to $\Lambda(\nu)$ proves the matching bound.

The input \eqref{eq:clusterinput} implies $\gamma_{600}=1$.
A given consecutive pair belongs to at most $601$ integer-start
windows of width $600$, so the associated cumulative gap count
has a positive-power lower bound along large cutoffs.
The trivial bound $G_H(X)\le X$ gives the reverse inequality
on its exponent. Hence $\Lambda(\nu)\ge1/(600+\nu)>0$.
Together with the fixed-rank proof this completes
Theorem~\ref{thm:spectrum}.

For the canonical upper bound use the minimum eventual gap,
two, and $\nu=1$ in \eqref{eq:cutoff}. Reducing $\varepsilon$
slightly absorbs the rounding, leaving all ascents below
$\EE((1/3+\varepsilon)k)$. Their number bounds $N_\eta(k)$.
Together with Corollary~\ref{cor:power}, this proves
Theorem~\ref{thm:canonical}.

\section{A finite phase profile}\label{sec:phase}
\begin{theorem}\label{thm:phase}
Let $H_*$ be the least $H\ge2$ for which $\gamma_H=1$.
This integer exists, and for every fixed $\nu>0$,
\[
 \Lambda(\nu)=\max_{2\le H\le H_*}\frac{\gamma_H}{H+\nu}.
\]
For all sufficiently large $\nu$,
\begin{equation}\label{eq:recovery}
 \Lambda(\nu)=\frac1{H_*+\nu},\qquad
 \frac1{\Lambda(\nu)}-\nu=H_*.
\end{equation}
\end{theorem}
\begin{proof}
Existence follows from $\gamma_{600}=1$.
If $H>H_*$, then
$\gamma_H/(H+\nu)\le1/(H+\nu)<1/(H_*+\nu)$.
For each $H<H_*$ put $\gamma=\gamma_H<1$.
With positive denominators, the inequality
$\gamma/(H+\nu)<1/(H_*+\nu)$ is equivalent to
$(1-\gamma)\nu>\gamma H_*-H$.
It holds beyond a finite threshold. Take the greatest threshold
over the finitely many smaller classes.
\end{proof}
Here full frequency means only $\gamma_H=1$. It need not imply
positive relative density or a fixed power lower bound at every
large cutoff. The theorem identifies a finite maximum without
determining its entries numerically.

\section{Arithmetic realization by affine prime sums}\label{sec:affine}
We prove Theorem~\ref{thm:affine} without discarding exceptional primes.

\subsection{Fixed-modulus limits and independence}
For fixed $Q$, the prime number theorem in arithmetic progressions gives
\[
 \frac{\#\{p\le X:p\text{ prime},\ p\equiv a\pmod Q\}}{\pi(X)}
       \longrightarrow\frac1{\varphi(Q)}\quad((a,Q)=1).
\]
Nonunit prime classes contain only primes dividing $Q$, so their
normalized contribution vanishes. Full Cartesian sampling of the
ordered pairs makes each fixed pair of unit classes have limiting
probability $\varphi(Q)^{-2}$.

Take $Q$ to be the product of any fixed finite set of primes.
The Chinese remainder theorem identifies the unit-pair space with
the product of the local unit-pair spaces. Divisibility of $F$
at a prime depends only on the local pair, so the limiting
divisibility indicators are independent, with exact law
\begin{equation}\label{eq:localcount}
 \eta_\ell=
 \frac{\#\{(u,v)\in(\mathbb F_\ell^\times)^2:
             \prod_j(c_j(u+v)+d_j)=0\}}{(\ell-1)^2}.
\end{equation}
For fixed $k,\ell$ the rank event uses only the primes at most
$\ell$. Summing their joint atoms gives \eqref{eq:density}.

If $F(p,q)=0$, then some $p+q=-d_j/c_j\le|d_j|
\le\sum_j|d_j|$. Only finitely many such pairs exist, independently
of $X$. Their normalized contribution is zero. The diagonal
was included throughout; its $\pi(X)$ pairs are also negligible
against $\pi(X)^2$. No modulus or rank was allowed to grow with $X$.

\subsection{Exact local counts}
Set $T_\ell=\{t\in\mathbb F_\ell:\prod_j(c_jt+d_j)=0\}$.
There are $\ell-1$ unit pairs with sum zero, and $\ell-2$
with any specified nonzero sum. Thus at every prime,
\begin{equation}\label{eq:rootcount}
 \eta_\ell=\frac{(\ell-2)|T_\ell|+
                    {\bf1}_{0\in T_\ell}}{(\ell-1)^2}.
\end{equation}
If a factor is identically zero modulo $\ell$,
$T_\ell=\mathbb F_\ell$ and $\eta_\ell=1$.
An empty root set gives zero, and coincident roots are counted
once. The formula covers all coefficient, shift, determinant,
and small-prime exceptions.

Outside the finite set dividing the nonzero coefficients,
nonzero shifts, or determinants $c_id_j-c_jd_i$, there are
$m$ distinct roots. At most one shift is zero; let $z_0$ be
one if it is and zero otherwise. Then
\begin{equation}\label{eq:genericlocal}
 \eta_\ell=\frac{m(\ell-2)+z_0}{(\ell-1)^2},\qquad
 \eta_\ell^{-1}=\frac{\ell}{m}-\frac{z_0}{m^2}+O_F(\ell^{-1}).
\end{equation}
The law satisfies \eqref{eq:smooth} with $\nu=m$.
Theorem~\ref{thm:spectrum} proves the arithmetic theorem.

The simplest instances illustrate the exceptional coordinate:
\[
\begin{array}{c|c|c}
 F(p,q)&\eta_2&\eta_\ell\quad(\ell>2)\\ \hline
 p+q&1&1/(\ell-1)\\
 p+q+1&0&(\ell-2)/(\ell-1)^2\\
 2(p+q)-1&0&(\ell-2)/(\ell-1)^2 .
\end{array}
\]
The forced factor $2$ in $p+q$ shifts subsequent ranks by one.
Both odd examples have exactly the canonical odd law.
All three have rate $\Lambda(1)$, despite different finite-rank
behavior. Equations \eqref{eq:rate} and \eqref{eq:recovery}
also show how sufficiently large fixed numbers of distinct affine
factors recover $H_*$. There is no uniformity claim when the family
varies with the rank.

\section{Late ascents and bounded prime gaps}\label{sec:late}
For a canonical law, a \emph{late ascent at coefficient $b$}
is an ascent across consecutive primes $p<q$ with
$\LL(p-1)\ge bk$. Unbounded supply requires existence beyond
every prescribed lower bound on both $k$ and $p$.

\begin{theorem}[Recurrence equivalence]\label{thm:late}
Let $H\ge1$ be an integer and
$1/(H+2)<b<1/(H+1)$. Gaps at most $H$ recur arbitrarily far out
if and only if for every $K,Q$ there exist $k\ge K$ and consecutive
primes $p<q$, with $p\ge Q$ and $\LL(p-1)\ge bk$, at which both
canonical densities ascend.
\end{theorem}
\begin{proof}
Choose $b<a<1/(H+1)$. To a sufficiently late gap at most $H$
assign $k=\lfloor\LL(p-1)/a\rfloor$. Then $k\to\infty$ with $p$,
$\LL(p-1)\ge ak\ge bk$, and $\LL(p-1)/k\to a$.
A small fixed band around $a$ in Lemma~\ref{lem:bands} gives
both ascents.

Conversely, $b>1/(H+2)$ and the one-sided ratio bound exclude
gaps at least $H+1$ from sufficiently late ascents. In the odd
law the correction to the ordinary threshold has absolute value
at most $1/(p-2)$, absorbed by the fixed strict margin.
Unbounded lower primes therefore give recurring gaps at most $H$.
\end{proof}

\begin{corollary}[Critical location]\label{cor:critical}
Suppose gaps at most $h\ge1$ recur arbitrarily far out and every
sufficiently late consecutive gap is at least $h$. The critical
canonical late-ascent coefficient is $1/(h+1)$: every smaller
positive coefficient has unbounded simultaneous ascent supply;
every larger coefficient forces both densities to descend for
all sufficiently large ranks and lower primes in that region.
No assertion at equality is made.
\end{corollary}
\begin{proof}
Use the forward rank assignment on gaps at most $h$ and the
one-sided upper ratio bound on gaps at least $h$.
\end{proof}

For every fixed $b>1/5$, an ascent in either canonical law at
sufficiently large $k,p$ with $\LL(p-1)\ge bk$ must cross a gap
strictly below four. Gaps between odd primes are positive even
integers, so it must be two. An independently proved unbounded
supply would therefore imply infinitely many twin primes.
That supply is not established here. A large finite ascent
certifies a finite prime pair, not recurrence.

\section{Implications and scope}\label{sec:implications}
The spectrum gives a quantitative dictionary. Bounded-gap counts
yield reversal counts through the finite selection lemma; the
upper comparison shows that no other late ascents can improve
the rate. Finite changes in an admissible law preserve that rate.
Many different affine prime-pair expressions consequently share
a spectrum despite different forced primes and small ranks.

Location and frequency encode different information.
Corollary~\ref{cor:critical} depends on the least recurring gap.
The rate $\Lambda(\nu)$ depends on how frequently bounded gap
classes occur. Infinitude alone need not give $\gamma_H>0$:
a hypothetical count growing as a power of $\log X$ has
$\gamma_H=0$. This is why bare infinitude cannot replace
the input to Corollary~\ref{cor:power}.

For Goldbach-type expressions, the arithmetic theorem describes
factorization in the full representation-weighted ensemble.
It gives no lower bound for representations of each fixed even
integer. Fixed-modulus independence likewise does not supply
uniform estimates when a prime cutoff grows with input size.
Moving-tail distribution laws require additional uniformity.

Repeated rises coexist with a final decreasing tail at each
rank. This is a statement about a probability sequence's shape,
rather than cancellation in a prime-counting error term.
No implication for the Riemann hypothesis, Robin's inequality,
or colossally abundant numbers is proved. Such an application
would require a new quantitative bridge to the corresponding
error term or optimization. The completed result is the spectrum
and its precise prime-gap transfers. Current exploratory
approaches to independent ascent supply are outside this paper.

\section{Formalization and reproducibility}\label{sec:formal}
The principal maintained owners are:
\begin{itemize}
\item \texttt{Assembly/ReciprocalSmoothSpectrum.lean}:
all-rank finite maxima and Theorem~\ref{thm:spectrum};
\item \texttt{Assembly/Headline.lean}: Theorem~\ref{thm:canonical};
\item \texttt{Assembly/AffineSpectrum.lean}:
Theorem~\ref{thm:affine} and arithmetic phase recovery;
\item \texttt{Assembly/GapSpectrumPhase.lean}:
Theorem~\ref{thm:phase};
\item \path{Assembly/GapRecurrence.lean} and\\
\path{Proof/Upper/GapScaleCutoff.lean}:
Theorem~\ref{thm:late}, Corollary~\ref{cor:critical}, and the
conditional inverse twin-prime transfer.
\end{itemize}
Paths are relative to \texttt{PrimeFactorOscillations/}. The
\href{https://github.com/kimihiro64/prime-factor-oscillations/blob/main/paper/THEOREM_STATUS.md}{theorem ledger}
maps intermediate arguments to declarations. Ordinary prose
sometimes combines several compiled lemmas; these combinations
are distinguished from individually named declarations.

The analytic inputs are the reciprocal-prime estimate,
fixed-modulus prime progression limits and \eqref{eq:clusterinput}.
Narrow leaves of the pinned Erd\H{o}s 690 development supply
ordinary density and symmetric bounds; its conditional
classification is not an input. The width-$600$ count is
extracted from a separately pinned prime-cluster development.
This paper proves the selection, transfers, spectrum and arithmetic
realization; it does not reproduce the underlying analytic
prime-distribution proofs.

The public
\href{https://github.com/kimihiro64/prime-factor-oscillations/blob/main/Challenge.lean}{Challenge}
states two headline theorems over Mathlib primitives.
\href{https://github.com/kimihiro64/prime-factor-oscillations/blob/main/Solution.lean}{Solution}
contains their actual proofs. The two deliberate Challenge
placeholders are absent from the Solution closures.
The recorded audit checks the public compiled types and twelve
theorem closures, finding only \texttt{propext},
\texttt{Classical.choice}, and \texttt{Quot.sound}.
An explicit hypothesis in a conditional theorem, such as ascent
supply, remains a hypothesis; an axiom audit does not remove it.

The
\href{https://github.com/kimihiro64/prime-factor-oscillations/blob/main/DEPENDENCY_REUSE.md}{dependency record}
and
\href{https://github.com/kimihiro64/prime-factor-oscillations/blob/main/BUILD.md}{build guide}
specify Windows Lean 4.34.0, source revisions and artifact hashes.
The project reuses the Erd\H{o}s dependency's published cache
without rebuilding or redistributing it. Additional analytic
compatibility artifacts still require an existing matching view;
a complete fresh-machine bootstrap is not claimed. The
\href{https://github.com/kimihiro64/prime-factor-oscillations/blob/main/paper/VERIFICATION.md}{verification record}
separates initial proof checks from this editorial revision.
Branch links provide mutable navigation; the records retain
checked source hashes.

No numerical experiment proves an asymptotic step here.
The exact inequality $1/602<1/601$ supplies the canonical
constant's strict margin; no effective onset rank is given.
Official Comparator/NanoDa replay and independent expert review
have not been claimed. Kernel checking verifies proof terms
for the stated definitions under the reported foundations.
Source fidelity, attribution, novelty and significance require
scrutiny beyond the kernel. AI assistance was used for research,
proof development and exposition.

\begin{thebibliography}{9}
\bibitem{Erdos}
P. Erd\H{o}s, \emph{Some unconventional problems in number theory},
Ast\'erisque \textbf{61} (1979), 73--82, especially pp.~74--75.
\href{https://www.numdam.org/item/AST_1979__61__73_0.pdf}{Original article}.

\bibitem{ET}
P. Erd\H{o}s and G. Tenenbaum,
\emph{Sur les densit\'es de certaines suites d'entiers},
Proc. London Math. Soc. (3) \textbf{59} (1989), 417--438,
Corollary~5 and equation~(1.20).
\href{https://doi.org/10.1112/plms/s3-59.3.417}{doi:10.1112/plms/s3-59.3.417}.

\bibitem{Cambie}
S. Cambie, \emph{Resolution of Erd\H{o}s' problems about unimodularity},
J. Number Theory \textbf{280} (2026), 271--277.
\href{https://doi.org/10.1016/j.jnt.2025.08.014}{doi:10.1016/j.jnt.2025.08.014};
\href{https://arxiv.org/abs/2501.10333v1}{arXiv:2501.10333v1}
(17 January 2025).

\bibitem{WC}
S. Wang and D. Crapis,
\emph{A Complete Answer to Erd\H{o}s Problem 690},
\href{https://arxiv.org/abs/2605.08542v1}{arXiv:2605.08542v1}
(8 May 2026).

\bibitem{Maynard}
J. Maynard, \emph{Dense clusters of primes in subsets},
Compositio Math. \textbf{152} (2016), 1517--1554.
\href{https://doi.org/10.1112/S0010437X16007296}{Journal version};
\href{https://arxiv.org/abs/1405.2593v2}{arXiv:1405.2593v2}.
\end{thebibliography}

\bigskip
\noindent Copyright 2026 Jonas Whidden. Original text and rendered paper:
Apache License 2.0 or CC BY 4.0, at the recipient's option.
\end{document}
